\clearpage
\VerdierPhysicalPageAnchor{0061}%
\VerdierSemanticAnchor{ch1-1-6-2}
\VerdierSemanticAnchor{ch1-1-6-2-proposition}
\noindent\textbf{命题 1.6.2.}\quad

\VerdierSemanticAnchor{ch1-1-6-2-part1}
\noindent\emph{1) 设
$\alpha\in Z^2(\mathbb Z^n,\mathbb Z/2\mathbb Z)$ 为正规化非齐次余循环，
$\dot\alpha\in H^2(\mathbb Z^n,\mathbb Z/2\mathbb Z)$ 为其上同调类，
$\dot\alpha_{i,j}\in H^2(\mathbb Z^n_{i,j},\mathbb Z/2\mathbb Z)$ 为其
各分量。则（利用 (1.6.1.2)）：}

\VerdierSemanticAnchor{ch1-1-6-2-eq1}
\begin{equation}
  \dot\alpha_{i,j}
  =\alpha(e_i,e_j)^{-1}\alpha(e_j,e_i).
  \tag{1.6.2.1}
\end{equation}

\VerdierSemanticAnchor{ch1-1-6-2-convention}
\noindent\emph{（$\mathbb Z/2\mathbb Z$ 的群运算采用乘法记号。）}

\VerdierSemanticAnchor{ch1-1-6-2-part2}
\noindent\emph{2) 设 $\alpha_1$ 与 $\alpha_2$ 是
$Z^2(\mathbb Z^n,\mathbb Z/2\mathbb Z)$ 中两个上同调等价的元素，并对每个
$i\in[1,n]$ 给定元素 $\beta_i\in\mathbb Z/2\mathbb Z$。则存在唯一映射}

\VerdierSemanticAnchor{ch1-1-6-2-beta-map}
\begin{equation*}
  \beta\colon\mathbb Z^n\longrightarrow\mathbb Z/2\mathbb Z
\end{equation*}
\noindent\emph{满足：}

\VerdierSemanticAnchor{ch1-1-6-2-eq2}
\begin{equation}
  \begin{aligned}
    \alpha_2(\sigma_1,\sigma_2)\alpha_1^{-1}(\sigma_1,\sigma_2)
      &=\beta(\sigma_2)\beta(\sigma_1+\sigma_2)^{-1}\beta(\sigma_1),\\[-1mm]
    &\hspace{7em}\text{其中 }\sigma_1,\sigma_2\in\mathbb Z^n,
  \end{aligned}
  \tag{1.6.2.2}
\end{equation}

\VerdierSemanticAnchor{ch1-1-6-2-eq3}
\begin{equation}
  \beta(e_i)=\beta_i.
  \tag{1.6.2.3}
\end{equation}

\VerdierSemanticAnchor{ch1-1-6-2-proof}
\VerdierSemanticAnchor{ch1-1-6-2-proof-first}
为证明第一个断言，我们归约到 $n=2$ 的情形，并直接验证当 $\alpha$ 加上一个
余边界时，量 $\alpha(e_i,e_j)^{-1}\alpha(e_j,e_i)$ 不变。当 $\alpha$ 是
余边界时，它是对称的，因此公式 (1.6.2.1) 在此情形成立。另一方面，余循环

\VerdierSemanticAnchor{ch1-1-6-2-alpha0}
\begin{equation*}
  \alpha_0(p_1e_i+q_1e_j,p_2e_i+q_2e_j)=(-1)^{q_1p_2}
\end{equation*}
不与 $0$ 上同调等价，因为

\VerdierSemanticAnchor{ch1-1-6-2-alpha0-test}
\begin{equation*}
  \alpha_0(e_i,e_j)^{-1}\alpha_0(e_j,e_i)=-1,
\end{equation*}
并且每个不与零上同调等价的余循环都与 $\alpha_0$ 上同调等价。

\VerdierSemanticAnchor{ch1-1-6-2-proof-second}
为证明第二个断言，注意 (1.6.2.2) 的任意两个解之差是一个群同态。因此总能
唯一地用一个群同态修正 (1.6.2.2) 的一个解，使它在各 $e_i$ 上取预先指定的值。

\VerdierSemanticAnchor{ch1-1-6-3}
\VerdierSemanticAnchor{ch1-1-6-3-definition}
\noindent\textbf{定义 1.6.3.}\quad
设 $e_1,\ldots,e_n$ 为 $\mathbb Z^n$ 的一组基。\emph{相对于基
$e_1,\ldots,e_n$ 的 Koszul 类}，是
$H^2(\mathbb Z^n,\mathbb Z/2\mathbb Z)$ 中唯一一个在各群
$H^2(\mathbb Z^n_{i,j},\mathbb Z/2\mathbb Z)$ 上的限制均非零的元素。

\VerdierSemanticAnchor{ch1-1-6-3-basis-dependence}
注意，Koszul 类的定义涉及 $\mathbb Z^n$ 的一组基的选择。一般而言，与两组
不同基相联系的 Koszul 类也不同。

\VerdierSemanticAnchor{ch1-1-6-3-filtration}
设 $e_1,\ldots,e_n$ 为一组按照 $[1,n]$ 上自然次序排列的基。以
$\mathbb Z^n(\leq i)$（分别以 $\mathbb Z^n(\geq i)$）表示由满足
$j\leq i$（分别满足 $j\geq i$）的 $e_j$ 生成的 $\mathbb Z^n$ 的子群。
